\subsection{MINUSCULE WEIGHTS}

PROPOSITION 6. Let \(\lambda \in \mathrm{P}\) , and \(\mathcal{X}\) the smallest R-saturated subset of P containing \(\lambda\) . Choose a norm \(\| \cdot \|\) as in Prop. 5. The following conditions are equivalent:

\begin{enumerate}
\def\labelenumi{(\roman{enumi})}
\item
  \(\mathcal{X} = \mathrm{W}.\lambda\) ;
\item
  all the elements of \(\mathcal{X}\) have the same norm;
\item
  for all \(\alpha \in \mathrm{R}\) , we have \(\lambda(H_{\alpha}) \in \{0,1,-1\}\) .
\end{enumerate}

Every non-empty R-saturated subset of P contains an element \(\lambda\) satisfying the above conditions.

Introduce the condition:

(ii') for all \(\alpha \in \mathrm{R}\) and for every integer \(t\) between 0 and \(\lambda(H_{\alpha})\) ,

\[
\| \lambda - t \alpha \| \geq \| \lambda \|.
\]

\begin{enumerate}
\def\labelenumi{(\roman{enumi})}
\tightlist
\item
  \(\Longrightarrow\) (ii) \(\Longrightarrow\) (ii'): This is clear.
\end{enumerate}

\((\mathrm{ii}^{\prime}) \Longrightarrow (\mathrm{iii})\) : Assume that condition \((\mathrm{ii}^{\prime})\) is satisfied. Let \(\alpha \in \mathrm{R}\) . We have \(\| \lambda \| = \| \lambda - \lambda (H_{\alpha})\alpha \|\) , so \(\| \lambda - t\alpha \| < \| \lambda \|\) for every integer \(t\) strictly between 0 and \(\lambda (H_{\alpha})\) ; hence, there can be no such integers, so \(|\lambda (H_{\alpha})| \leq 1\) .

\begin{enumerate}
\def\labelenumi{(\roman{enumi})}
\setcounter{enumi}{2}
\tightlist
\item
  \(\Longrightarrow\) (i): Assume that condition (iii) is satisfied. Let \(w \in W\) and \(\alpha \in R\) . Then \((w\lambda)(H_{\alpha}) = \lambda(H_{w^{-1}\alpha}) \in \{0,1,-1\}\) ; thus, if t is an integer between 0 and \((w\lambda)(H_{\alpha})\) , \(w\lambda - t\alpha\) is equal to \(w\lambda\) or \(s_{\alpha}(w\lambda)\) . This proves that W. \(\lambda\) is R-saturated, so X = W. \(\lambda\) .
\end{enumerate}

Let Y be a non-empty R-saturated subset of P. There exists in Y an element \(\lambda\) of minimum norm. It is clear that \(\lambda\) satisfies condition (ii'), hence the last assertion of the proposition.

PROPOSITION 7. Let V be a finite dimensional simple g-module, X the set of weights of V, and \(\lambda\) the highest element of X (cf.~Prop. 5 (i)). Conditions (i), (ii) and (iii) of Prop. 6 are equivalent to:

\begin{enumerate}
\def\labelenumi{(\roman{enumi})}
\setcounter{enumi}{3}
\tightlist
\item
  for all \(\alpha \in \mathrm{R}\) and all \(x \in \mathfrak{g}^{\alpha}\) , we have \((x_{\mathrm{V}})^2 = 0\) .
\end{enumerate}

If these conditions are satisfied, all the weights of V have multiplicity 1.

If (i) is satisfied, then \(X = W \cdot \lambda\) and the weights all have the same multiplicity as \(\lambda\) (Cor. 2 of Prop. 2), in other words, multiplicity 1. Moreover, if \(w \in W\) and \(\alpha \in R\) , \(w\lambda + t\alpha\) cannot be a weight of V unless \(|t| \leq 1\) ; thus, if \(x \in g^{\alpha}\) ,

\[
(x _ {\mathrm{V}}) ^ {2} (\mathrm{V} ^ {w (\lambda)}) \subset \mathrm{V} ^ {w (\lambda) + 2 \alpha} = 0,
\]

so \((x_{\mathrm{V}})^2 = 0\) , which proves that (i) \(\Longrightarrow\) (iv).

Conversely, assume that (iv) is satisfied. Let \(\alpha \in \mathrm{R}\) , and give V the \(\mathfrak{sl}(2,k)\) -module structure defined by the elements \(X_{\alpha}, X_{-\alpha}, H_{\alpha}\) of \(\mathfrak{g}\) . Condition (iv), applied to \(x = X_{\alpha}\) , implies that the weights of the \(\mathfrak{sl}(2,k)\) -module V belong to \(\{0,1,-1\}\) (cf.~§1, no. 2, Cor. of Prop. 2). In particular, \(\lambda(H_{\alpha}) \in \{0,1,-1\}\) , so (iv) \(\Longrightarrow\) (iii).

PROPOSITION 8. Assume that g is simple. Denote by \(\alpha_{1},\ldots,\alpha_{l}\) the elements of B. Let \(\varpi_{1},\ldots,\varpi_{l}\) be the corresponding fundamental weights. Let \(H = n_{1}H_{\alpha_{1}} + \cdots + n_{l}H_{\alpha_{l}}\) be the highest root of \(R^{\vee}\) , and J the set of \(i \in \{1,\ldots,l\}\) such that \(n_{i} = 1\) . Let \(\lambda \in P_{++} - \{0\}\) . Then conditions (i), (ii) and (iii) of Prop. 6 are equivalent to each of the following conditions:

\begin{enumerate}
\def\labelenumi{(\alph{enumi})}
\setcounter{enumi}{21}
\tightlist
\item
  \(\lambda(H)=1;\)
\end{enumerate}

\begin{enumerate}
\def\labelenumi{(\roman{enumi})}
\setcounter{enumi}{5}
\tightlist
\item
  there exists \(i \in \mathrm{J}\) such that \(\lambda = \varpi_i\) .
\end{enumerate}

The \(\varpi_{i}\) , for \(i \in J\) , form a system of representatives in P(R) of the non-zero elements of P(R)/Q(R).

Let \(\lambda = u_{1}\varpi_{1} + \cdots + u_{l}\varpi_{l}\) , where \(u_{1}, \ldots, u_{l}\) are integers \(\geq 0\) and not all zero. Then \(\lambda(H) = u_{1}n_{1} + \cdots + u_{l}n_{l}\) and \(n_{1} \geq 1, \ldots, n_{l} \geq 1\) , which gives the equivalence of (v) and (vi) immediately. On the other hand, \(\lambda(H) = \sup_{\alpha \in \mathrm{R}_{+}} \lambda(H_{\alpha})\) , and \(\lambda(H) > 0\) since \(\lambda\) is a non-zero element of \(P_{++}\) . Hence condition (v) is equivalent to the condition \(\lambda(H_{\alpha}) \in \{0, 1\}\) for all \(\alpha \in R\) , in other words to condition (iii) of Prop. 6.

The last assertion of the proposition follows from Chap. VI, §2, no. 3, Cor. of Prop. 6.

DEFINITION 1. Assume that \(\mathfrak{g}\) is simple. A minuscule weight of \((\mathfrak{g},\mathfrak{h})\) is an element of \(\mathrm{P}_{++} - \{0\}\) which satisfies the equivalent conditions (i), (ii), (iii), (iv), (v) and (vi) of Prop. 6, 7 and 8.

Remark. Assume that g is simple. Let \(\Sigma^{\prime\vee}\) be the Coxeter graph of the affine Weyl group \(\mathrm{W}_{a}(\mathrm{R}^{\vee})\) . Recall that the vertices of \(\Sigma^{\prime\vee}\) are the vertices of the Coxeter graph \(\Sigma^{\vee}\) of \(\mathrm{W}(\mathrm{R}^{\vee})\) , together with a supplementary vertex 0. The group \(\mathrm{A}(\mathrm{R}^{\vee})\) operates on \(\Sigma^{\prime\vee}\) leaving 0 fixed. The group \(\mathrm{Aut}(\Sigma^{\prime\vee})\) is canonically isomorphic to the semi-direct product of \(\mathrm{A}(\mathrm{R}^{\vee})/\mathrm{W}(\mathrm{R}^{\vee})\) with a group \(\Gamma_{C}\) (cf.~Chap. VI, §2, no. 3, and Chap. VI, §4, no. 3); clearly \((\operatorname{Aut} \Sigma'^{\vee})(0) = \Gamma_C(0)\) ; and \(\Gamma_C(0)\) consists of 0 and the vertices of \(\Sigma^{\vee}\) corresponding to the \(\varpi_i\) for \(i \in \mathbf{J}\) (cf.~Chap. VI, §2, Prop. 5 and Remark 1 of no. 3). In summary, the minuscule weights are the fundamental weights corresponding to the vertices of \(\Sigma^{\vee}\) which can be obtained from 0 by the operation of an element of \(\operatorname{Aut}(\Sigma'^{\vee})\) .

With the notations of Chap. VI, Plates I to IX, we deduce from the preceding that the minuscule weights are the following:

For type \(A_{l}\) ( \(l \geq 1\) ): \(\varpi_{1}, \ldots, \varpi_{l}\) .

For type \(B_{l}\) ( \(l \geq 2\) ): \(\varpi_{l}\) .

For type \(C_{l}\) ( \(l \geq 2\) ): \(\varpi_{1}\) .

For type \(D_{l}\) ( \(l \geq 3\) ): \(\varpi_{1}, \varpi_{l-1}, \varpi_{l}\) .

For type \(E_{6}\) : \(\varpi_{1}, \varpi_{6}\) .

For type \(\mathrm{E}_7\) : \(\varpi_7\) .

For types \(E_{8}, F_{4}, G_{2}\) there are no minuscule weights.

